\appendix
\part{Anhang}

\newcommand{\lit}[3]{\vspace{5pt}\noindent\textbf{[#1]} #2\par\nopagebreak\vspace{1pt}{\small\itshape\color{grau} #3}\par\vspace{2pt}}

\chapter{Formelsammlung und Übersicht}

\section{Begriffe und Formeln}

\begin{merksatz}
\boxtitel{Teil I und II}
\begin{center}
{\small
\renewcommand{\arraystretch}{1.6}
\begin{tabular}{@{}P{4.4cm}P{4.2cm}@{}}
\toprule
$Y_i(1),\ Y_i(0)$ & potentielle Ergebnisse (2.1)\\
$Y_i = Y_i(X_i)$ & Konsistenz\\
$\tau = \mathbb{E}[Y(1)-Y(0)]$ & Durchschnittseffekt (2.3)\\
$\tau_{\text{beh}}$ & Effekt bei den Behandelten\\
$(Y(1),Y(0)) \perp X \mid Z$ & bedingte Unabhängigkeit\\
$0 < \Pr[X{=}1\mid Z] < 1$ & Positivität\\
$\sum_z \mathbb{E}[Y\mid x,z]\Pr[z]$ & Anpassungsformel (2.4)\\
$\Pr[Y\mid\dops(x)]$ & Eingriffsverteilung (1.1)\\
$n = \frac{(z_{1-\alpha/2}+z_{1-\beta})^{2}[\,\cdot\,]}{(p_0-p_1)^{2}}$ & Fallzahl (3)\\
$\prod_i \Pr[v_i \mid \text{Eltern}]$ & Faktorisierung (4.2)\\
Kette, Gabel, Kollider & drei Bausteine (4)\\
d-Trennung & alle Wege gesperrt (4.3)\\
kein Nachkomme, Hinterwege zu & Hintertür (6.1)\\
$\sum_m \Pr[m\mid x]\sum_{x'}\Pr[Y\mid m,x']\Pr[x']$ & Vordertür (6.3)\\
\bottomrule
\end{tabular}}
\end{center}
\end{merksatz}

\begin{merksatz}
\boxtitel{Teil III bis V}
\begin{center}
{\small
\renewcommand{\arraystretch}{1.6}
\begin{tabular}{@{}P{4.4cm}P{4.2cm}@{}}
\toprule
$e(\vek z) = \Pr[X{=}1\mid\vek z]$ & Behandlungsneigung (7.1)\\
$\perp X \mid e(Z)$ & Rosenbaum-Rubin (7.2)\\
$w = 1/e$ bzw.\ $1/(1-e)$ & Gewichte (8.1)\\
$\mathbb{E}[XY/e(Z)] = \mathbb{E}[Y(1)]$ & (8.2)\\
$\frac{X(Y-\hat m_1)}{\hat e} + \hat m_1$ & doppelt abgesichert (8.3)\\
$(\sum w)^{2}/\sum w^{2}$ & wirksame Fallzahl\\
$\frac{\Delta\mathbb{E}[Y]}{\Delta\mathbb{E}[X]}$ & Wald-Schätzer (9.1)\\
$F > 10$ & Instrumentstärke\\
Sprung an der Schwelle & Regressionssprung (10)\\
$(\bar y_{A,1}-\bar y_{A,0}) - (\bar y_{B,1}-\bar y_{B,0})$ & doppelte Differenz (10)\\
sieben Zeilen & Zielstudie (11)\\
$\sum_\ell \Pr[Y\mid a_0,\ell,a_1]\Pr[\ell\mid a_0]$ & g-Formel (12.1)\\
$\prod_t 1/\Pr[A_t \mid \text{Geschichte}]$ & Gewichte über die Zeit\\
$c' + ab$ & Mediation, linear (13)\\
$\tau(\vek z)$ & bedingter Effekt (14)\\
\bottomrule
\end{tabular}}
\end{center}
\end{merksatz}

\section{Die Landkarte}

\begin{merksatz}
\boxtitel{Welcher Satz worauf beruht}
\begin{center}
{\small
\renewcommand{\arraystretch}{1.4}
\begin{tabular}{@{}P{3.8cm}P{1.0cm}P{3.4cm}@{}}
\toprule
\textbf{Ergebnis} & \textbf{Kap.} & \textbf{beruht auf}\\
\midrule
Beobachtungs\-äquivalenz & $1$ & zwei Modelle, eine Tabelle\\
kein individueller Effekt & $2$ & Begriffsbildung\\
Anpassungsformel & $2$ & totale Erwartung\\
Randomisierung & $3$ & Unabhängigkeit vor der Zuteilung\\
Faktorisierung & $4$ & Modellaufbau\\
d-Trennung & $4$ & Wege und Kollider\\
Hintertür & $6$ & Durchtrennen der Pfeile\\
Vordertür & $6$ & zwei Stücke, je unverzerrt\\
Neigung genügt & $7$ & Turmregel\\
Gewichtung & $8$ & Konsistenz und Kürzen\\
doppelte Absicherung & $8$ & ein Faktor wird null\\
Wald-Schätzer & $9$ & Ausschlussbedingung\\
Befolgereffekt & $9$ & vier Patiententypen\\
g-Formel & $12$ & Zeitordnung\\
Mediation & $13$ & zwei Welten zugleich\\
\bottomrule
\end{tabular}}
\end{center}
\end{merksatz}

\section{Was nicht bewiesen wurde}

\begin{merksatz}
\boxtitel{Vier Auslassungen}
\begin{center}
{\small
\renewcommand{\arraystretch}{1.5}
\begin{tabular}{@{}P{3.4cm}P{1.2cm}P{3.4cm}@{}}
\toprule
\textbf{Ergebnis} & \textbf{Kap.} & \textbf{Wo es steht}\\
\midrule
d-Trennungssatz & $4$ & Pearl\zit{18}, Kap.\ 1\\
Hintertürformel & $6$ & Pearl\zit{18}, Kap.\ 3\\
Vollständigkeit des do-Kalküls & $6$ & Pearl u.\,a.\zit{19}\\
Satz über natürliche Effekte & $13$ & VanderWeele\zit{26}\\
\bottomrule
\end{tabular}}
\end{center}

\begin{center}
\begin{minipage}{0.88\textwidth}
\kursiv{\textbf{Alle übrigen Sätze des Bandes sind vollständig bewiesen} --- darunter die Anpassungsformel, Satz 3.1 über die Randomisierung, Satz 7.2 über die Behandlungsneigung, Satz 8.2 über die Gewichtung und Satz 8.3 über die doppelte Absicherung.}
\end{minipage}
\end{center}
\end{merksatz}

\section{Zeittafel}

\begin{merksatz}
\boxtitel{Von Hume bis zur Zielstudie}
\begin{center}
{\small
\renewcommand{\arraystretch}{1.32}
\begin{tabular}{@{}P{1.3cm}P{3.0cm}P{4.0cm}@{}}
\toprule
\textbf{Jahr} & \textbf{Person} & \textbf{Ereignis}\\
\midrule
$1748$ & Hume & Ursache als beständige Verknüpfung\\
$1843$ & Mill & Methoden der Übereinstimmung\\
$1854$ & Snow & die Wasserwerke Londons\\
$1899$ & Pearson & die Umkehr in Teiltabellen\\
$1918$ & S.\ Wright & Pfadanalyse\\
$1923$ & Neyman & potentielle Ergebnisse\\
$1928$ & P.\ Wright & Instrumentvariablen\\
$1935$ & Fisher & \kursiv{The Design of Experiments}\\
$1946$ & Berkson & Verzerrung durch Aufnahme\\
$1948$ & Bradford Hill & Streptomycinstudie\\
$1951$ & Simpson & das Paradoxon\\
$1952$ & Horvitz, Thompson & Gewichtung nach Kehrwert\\
$1959$ & Cornfield & Empfindlichkeits\-argument\\
$1960$ & Thistlethwaite & Regressionssprung\\
$1965$ & Bradford Hill & neun Gesichtspunkte\\
$1974$ & Rubin & potentielle Ergebnisse für Beobachtungen\\
$1983$ & Rosenbaum, Rubin & Behandlungsneigung\\
$1986$ & \textbf{Robins} & \textbf{g-Formel}\\
$1986$ & Holland & \enquote{keine Ursache ohne Manipulation}\\
$1986$ & Baron, Kenny & Mediation in drei Regressionen\\
$1993$ & \textbf{Pearl} & \textbf{Hintertürkriterium}\\
$1994$ & Imbens, Angrist & Befolgereffekt\\
$1994$ & Robins u.\,a. & doppelte Absicherung\\
$1995$ & Pearl & Vordertürformel\\
$2000$ & Pearl & \kursiv{Causality}\\
$2002$ & WHI & Umkehr des Hormonbefundes\\
$2006$ & Shpitser, Pearl & Vollständigkeit des do-Kalküls\\
$2008$ & Hernán u.\,a. & Auflösung des Hormonwiderspruchs\\
$2016$ & Hernán, Robins & Zielstudie\\
$2018$ & Athey, Wager & kausale Wälder\\
$2021$ & Card, Angrist, Imbens & Nobelpreis\\
\bottomrule
\end{tabular}}
\end{center}
\end{merksatz}

\section{Glossar}

\begin{merksatz}
\boxtitel{Begriffe von A bis Z}
\begin{center}
{\small
\renewcommand{\arraystretch}{1.3}
\begin{longtable}{@{}P{3.6cm}P{5.0cm}@{}}
\toprule
\textbf{Begriff} & \textbf{Bedeutung}\\
\midrule
\endhead
Behandlungsabsicht & Auswertung nach Zuteilung (3)\\
Behandlungsneigung & $\Pr[X{=}1\mid Z]$ (7.1)\\
Befolger & wer der Zuteilung folgt (9)\\
d-Trennung & alle Wege gesperrt (4.3)\\
Doppelte Absicherung & eines von zwei Modellen genügt (8.3)\\
Doppelte Differenz & Trend und Ort heraus (10)\\
Gabel & gemeinsame Ursache (4)\\
g-Formel & Verlauf erzeugen, dann mitteln (12.1)\\
Hinterweg & Pfeil in $X$ hinein (6)\\
Indikationsverzerrung & Kranke werden anders behandelt (1)\\
Instrument & Ursache nur von $X$ (9)\\
Kollider & gemeinsame Folge (4)\\
Konsistenz & $Y = Y(X)$ (2)\\
Kontrafaktisch & was gewesen wäre (2.1)\\
Kreuzwelt & zwei Welten zugleich (13)\\
Mediator & auf dem Weg $X \to Y$ (5)\\
Positivität & beide Behandlungen kommen vor (2)\\
Randstrukturmodell & Modell der potentiellen Ergebnisse (12)\\
Regressionssprung & Schwelle als Zufall (10)\\
Störgröße & Ursache von $X$ und $Y$ (5)\\
Triangulation & verschiedene Annahmen, gleiches Ergebnis (11)\\
Unsterbliche Zeit & Spanne ohne Sterbemöglichkeit (11)\\
Vordertür & Effekt über den Mediator (6.3)\\
Zielstudie & erst die Studie, dann die Daten (11)\\
\bottomrule
\end{longtable}}
\end{center}
\end{merksatz}
